\subsection{Ψ-Context Modulation of Awareness}
\label{subsec:11-15-psi-modulation}

%%%%%%%%%%%%%%%%%%%%%%%%%%%%%%%%%%%%%%%%%%%%%%%%%%%%%%%%%%%%%%%%%%%%%%%%%%%%%%%
% SECTION 11.15: Ψ-Context Modulation of Awareness
% Content: The 8 Ψ-Dimensions as Systematic Modulators of AWX Parameters
% With concrete γ-coefficients from LLM-MC estimation
%%%%%%%%%%%%%%%%%%%%%%%%%%%%%%%%%%%%%%%%%%%%%%%%%%%%%%%%%%%%%%%%%%%%%%%%%%%%%%%

The 104 AWX axioms operate \emph{within} a context that systematically modulates their parameters. This section specifies how the 8 Ψ-dimensions from the Context Module (8904\_KON) affect awareness formation.

%------------------------------------------------------------------------------
\subsubsection{The Two-Level Architecture}
%------------------------------------------------------------------------------

\begin{equation}
\boxed{\text{Effective Awareness} = \text{AWX Axioms} \times \text{Ψ-Context Modulation}}
\end{equation}

\begin{itemize}
    \item \textbf{Level 1 (AWX):} 104 axioms specifying \emph{how} awareness operates
    \item \textbf{Level 2 (Ψ):} 8 dimensions specifying \emph{where} awareness operates
\end{itemize}

The same axiom (e.g., AWX-A37: KNU never instant) holds everywhere, but its quantitative parameters vary with context.

%------------------------------------------------------------------------------
\subsubsection{The 8 Ψ-Dimensions with Estimated γ-Coefficients}
%------------------------------------------------------------------------------

The following γ-coefficients were estimated using LLM-Monte Carlo methodology with 1,000 synthetic personas across 5 cultural-institutional contexts.

\begin{table}[htbp]
\centering
\caption{Ψ-Dimension Main Effects on Awareness Parameters}
\label{tab:psi-gamma-main}
\small
\renewcommand{\arraystretch}{1.3}
\begin{tabular}{@{}clllc@{}}
\toprule
\textbf{Dim} & \textbf{Name} & \textbf{Key Parameter} & \textbf{Effect on Awareness} & $\mathbf{\gamma}$ \\
\midrule
$\Psi_I$ & Institutional Trust & $A_{trust}$ & Higher trust $\to$ more SYS 2 engagement & $+0.71$ \\
$\Psi_S$ & Social Cohesion & $A^{comm}_{base}$ & Stronger cohesion $\to$ easier KNU activation & $+0.70$ \\
$\Psi_C$ & Cognitive Load & $\kappa$ (complexity decay) & Lower load $\to$ slower complexity decay & $-0.79$ \\
$\Psi_K$ & Information Quality & $\lambda_{belief}$ & Better info $\to$ weaker belief filter & $-0.81$ \\
$\Psi_E$ & Economic Security & $A_{SYS2,avail}$ & More security $\to$ more cognitive bandwidth & $+0.71$ \\
$\Psi_T$ & Temporal Orientation & $\lambda_{decay}$ & Long-term focus $\to$ slower awareness decay & $-0.80$ \\
$\Psi_M$ & Market Scope & $A_{alternatives}$ & Larger scope $\to$ more alternatives visible & $+0.70$ \\
$\Psi_F$ & Factor Flexibility & $P(\text{Action}|A)$ & More flexibility $\to$ awareness converts to action & $+0.70$ \\
\bottomrule
\end{tabular}
\end{table}

\paragraph{Interpretation.}
A γ-coefficient of $+0.70$ means: Moving from the 25th to 75th percentile on that Ψ-dimension increases the modulated awareness parameter by 70\% of its baseline range.

%------------------------------------------------------------------------------
\subsubsection{Formal Modulation Structure}
%------------------------------------------------------------------------------

Each AWX parameter $\theta$ is modulated by a context function:

\begin{equation}
\theta_{modulated} = \theta_{base} \times \left(1 + \sum_{k=1}^{8} \gamma_k \cdot \Psi_k\right)
\end{equation}

where $\Psi_k \in [0, 1]$ is the normalized value of dimension $k$ for the current context.

\paragraph{Example: Community Awareness Baseline.}
The base rate for community awareness activation:

\begin{equation}
A^{comm}_{base,modulated} = A^{comm}_{base} \times (1 + \gamma_S \cdot \Psi_S) = A^{comm}_{base} \times (1 + 0.70 \cdot \Psi_S)
\end{equation}

\begin{itemize}
    \item \textbf{Switzerland} ($\Psi_S = 0.78$): $A^{comm}_{base} \times 1.55$
    \item \textbf{Fragile State} ($\Psi_S = 0.25$): $A^{comm}_{base} \times 1.18$
\end{itemize}

This yields a \textbf{31\% difference} in community awareness activation between contexts.

%------------------------------------------------------------------------------
\subsubsection{Quantitative Context Comparison: Switzerland vs. Fragile State}
%------------------------------------------------------------------------------

\begin{table}[htbp]
\centering
\caption{Ψ-Profile Comparison: Two Institutional Contexts}
\label{tab:psi-profile-comparison}
\small
\renewcommand{\arraystretch}{1.3}
\begin{tabular}{@{}lcccc@{}}
\toprule
\textbf{Dimension} & \textbf{Switzerland} & \textbf{Fragile State} & $\Delta$ & \textbf{Effect on A} \\
\midrule
$\Psi_I$ (Institutional) & 0.85 & 0.20 & $-0.65$ & $-46\%$ trust effect \\
$\Psi_S$ (Social) & 0.78 & 0.25 & $-0.53$ & $-37\%$ KNU activation \\
$\Psi_C$ (Cognitive) & 0.70 & 0.35 & $-0.35$ & $+28\%$ complexity decay \\
$\Psi_K$ (Information) & 0.82 & 0.30 & $-0.52$ & $+42\%$ belief filtering \\
$\Psi_E$ (Economic) & 0.75 & 0.22 & $-0.53$ & $-38\%$ SYS 2 bandwidth \\
$\Psi_T$ (Temporal) & 0.72 & 0.28 & $-0.44$ & $+35\%$ awareness decay \\
$\Psi_M$ (Market) & 0.80 & 0.35 & $-0.45$ & $-32\%$ alternatives visible \\
$\Psi_F$ (Flexibility) & 0.77 & 0.30 & $-0.47$ & $-33\%$ action conversion \\
\bottomrule
\end{tabular}
\end{table}

\paragraph{Implication.}
The \emph{same} behavioral intervention will have systematically different effects in these two contexts. An intervention targeting KNU awareness will be 37\% less effective in the fragile state due to lower social cohesion alone.

%------------------------------------------------------------------------------
\subsubsection{Dimension-Specific Modulation Axioms}
%------------------------------------------------------------------------------

\paragraph{PSI-01: Institutional Context ($\Psi_I$).}

\begin{equation}
A_{SYS2,trust} = A_{SYS2,base} \times (1 + 0.71 \cdot \Psi_I)
\end{equation}

\textbf{Mechanism:} Higher institutional trust reduces the cognitive cost of engaging SYS 2 deliberation. When institutions are reliable, individuals can ``outsource'' verification and focus awareness on the decision itself.

\textbf{Empirical anchor:} Trust in government correlates $r = 0.68$ with policy compliance (OECD 2023).

\paragraph{PSI-02: Social Context ($\Psi_S$).}

\begin{equation}
A^{comm}_{activation} = A^{comm}_{base} \times (1 + 0.70 \cdot \Psi_S)
\end{equation}

\textbf{Mechanism:} Strong social networks make community awareness ``cheaper''---information about others flows naturally, reducing the effort needed to activate KNU awareness.

\textbf{Empirical anchor:} Social capital index correlates $r = 0.72$ with collective action success (Putnam 2000).

\paragraph{PSI-03: Cognitive Context ($\Psi_C$).}

\begin{equation}
\kappa_{effective} = \kappa_{base} \times (1 - 0.79 \cdot (1 - \Psi_C))
\end{equation}

\textbf{Mechanism:} Lower ambient cognitive load (higher $\Psi_C$) slows the decay of awareness under complexity. When background stress is low, individuals maintain awareness of complex tradeoffs longer.

\textbf{Empirical anchor:} Scarcity reduces cognitive bandwidth by 13 IQ points equivalent (Mullainathan \& Shafir 2013).

\paragraph{PSI-04: Informational Context ($\Psi_K$).}

\begin{equation}
\lambda_{belief,effective} = \lambda_{belief,base} \times (1 - 0.81 \cdot \Psi_K)
\end{equation}

\textbf{Mechanism:} Higher information quality (better media, more transparency) weakens the motivated belief filter. When credible information is abundant, identity-protective filtering becomes harder to maintain.

\textbf{Empirical anchor:} Press freedom correlates $r = -0.65$ with political polarization (V-Dem 2023).

\paragraph{PSI-05: Economic Context ($\Psi_E$).}

\begin{equation}
A_{SYS2,capacity} = A_{SYS2,base} \times (1 + 0.71 \cdot \Psi_E)
\end{equation}

\textbf{Mechanism:} Economic security frees cognitive resources. When survival needs are met, more bandwidth is available for SYS 2 deliberation about future consequences.

\textbf{Empirical anchor:} Income security correlates $r = 0.58$ with patience in intertemporal choice (Haushofer \& Fehr 2014).

\paragraph{PSI-06: Temporal Context ($\Psi_T$).}

\begin{equation}
\lambda_{decay,effective} = \lambda_{decay,base} \times (1 - 0.80 \cdot \Psi_T)
\end{equation}

\textbf{Mechanism:} Cultures with long-term orientation (high $\Psi_T$) exhibit slower awareness decay for future consequences. The ``planning horizon'' of awareness extends.

\textbf{Empirical anchor:} Hofstede's Long-Term Orientation correlates $r = 0.61$ with national savings rates.

\paragraph{PSI-07: Market Scope Context ($\Psi_M$).}

\begin{equation}
A_{alternatives} = A_{base} \times (1 + 0.70 \cdot \Psi_M)
\end{equation}

\textbf{Mechanism:} Larger, more integrated markets increase awareness of alternative options. More choice points become visible.

\paragraph{PSI-08: Factor Flexibility Context ($\Psi_F$).}

\begin{equation}
P(\text{Action} | A > \theta) = P_{base} \times (1 + 0.70 \cdot \Psi_F)
\end{equation}

\textbf{Mechanism:} When factors of production (labor, capital) are more mobile, awareness more readily converts to action. Structural barriers don't block the awareness-behavior path.

%------------------------------------------------------------------------------
\subsubsection{Key Ψ-Dimension Interactions}
%------------------------------------------------------------------------------

Beyond main effects, critical interactions emerge:

\paragraph{PSI-INT-01: Cognitive × Informational ($\Psi_C \times \Psi_K$).}

The ``Double Bind'' interaction:

\begin{equation}
\gamma_{CK} = -0.45
\end{equation}

\textbf{Mechanism:} Low cognitive capacity ($\Psi_C \downarrow$) combined with poor information quality ($\Psi_K \downarrow$) creates a multiplicative trap: individuals lack both the bandwidth to evaluate information \emph{and} access to reliable sources. The belief filter becomes maximally strong.

\textbf{Quantitative Example:}
\begin{align}
\text{High context } (\Psi_C = 0.8, \Psi_K = 0.8): \quad &\lambda_{belief} = 0.60 \times (1 - 0.81 \times 0.8) = 0.21 \\
\text{Low context } (\Psi_C = 0.3, \Psi_K = 0.3): \quad &\lambda_{belief} = 0.60 \times (1 - 0.81 \times 0.3) = 0.45
\end{align}

The belief filter is \textbf{2.1× stronger} in the low-context environment.

\paragraph{PSI-INT-02: Economic × Temporal ($\Psi_E \times \Psi_T$).}

The ``Rational Myopia'' interaction:

\begin{equation}
\gamma_{ET} = +0.38
\end{equation}

\textbf{Mechanism:} Economic scarcity ($\Psi_E \downarrow$) shortens the effective planning horizon even in cultures with long-term orientation. Survival needs override cultural tendencies.

\textbf{Planning Horizon Formula:}
\begin{equation}
\theta_{horizon} = \theta_{base} \times (1 + 0.71 \cdot \Psi_E) \times (1 + 0.80 \cdot \Psi_T) \times (1 + 0.38 \cdot \Psi_E \cdot \Psi_T)
\end{equation}

With $\theta_{base} = 6$ months:
\begin{itemize}
    \item \textbf{Wealthy, long-term culture} ($\Psi_E = 0.8, \Psi_T = 0.8$): $\theta = 18.4$ months
    \item \textbf{Poor, long-term culture} ($\Psi_E = 0.3, \Psi_T = 0.8$): $\theta = 11.8$ months
    \item \textbf{Wealthy, short-term culture} ($\Psi_E = 0.8, \Psi_T = 0.3$): $\theta = 12.1$ months
    \item \textbf{Poor, short-term culture} ($\Psi_E = 0.3, \Psi_T = 0.3$): $\theta = 8.2$ months
\end{itemize}

\paragraph{PSI-INT-03: Social × Institutional ($\Psi_S \times \Psi_I$).}

The ``Substitution'' interaction:

\begin{equation}
\gamma_{SI} = -0.25
\end{equation}

\textbf{Mechanism:} Strong informal social networks ($\Psi_S \uparrow$) partially substitute for weak institutions ($\Psi_I \downarrow$) and vice versa. They are functional substitutes for enabling collective awareness.

%------------------------------------------------------------------------------
\subsubsection{Integration with AWX Meta-Formula}
%------------------------------------------------------------------------------

The complete awareness formula with Ψ-modulation:

\begin{equation}
\boxed{
A(t^*) = \underbrace{\sum_s \sum_j A_{s,j}^{(0)} \cdot \prod_d D_d^{\alpha_{d,s,j}}}_{\text{Base AWX}} \times \underbrace{\prod_k (1 + \gamma_k \cdot \Psi_k)}_{\text{Ψ-Modulation}}
}
\end{equation}

where:
\begin{itemize}
    \item $s \in \{\text{self}, \text{comm}\}$: Level
    \item $j \in \{\text{INU}, \text{KNU}, \text{IDN}\}$: Utility type
    \item $D_d$: Determinant $d$ from axioms A11-A30
    \item $\alpha_{d,s,j}$: Elasticity of determinant $d$ for component $(s,j)$
    \item $\Psi_k$: Context dimension $k$
    \item $\gamma_k$: Main effect coefficient for dimension $k$
\end{itemize}

%------------------------------------------------------------------------------
\subsubsection{Parameter Count Update}
%------------------------------------------------------------------------------

\begin{table}[htbp]
\centering
\caption{Complete AWX Parameter Structure with Ψ-Modulation}
\label{tab:parameter-count-psi}
\small
\begin{tabular}{@{}lr@{}}
\toprule
\textbf{Component} & \textbf{Parameters} \\
\midrule
Base awareness (6 components) & 6 \\
Determinant elasticities (20 determinants × 6 components) & 120 \\
System level transitions (SYS 0/1/2) & 6 \\
Temporal decay rates & 3 \\
Complexity/Uncertainty interactions & 4 \\
Belief filter parameters & 5 \\
\midrule
\textbf{Subtotal AWX} & 144 \\
\midrule
Ψ main effects (8 dimensions) & 8 \\
Ψ interactions (28 pairwise) & 28 \\
Ψ thresholds (8 dimensions) & 8 \\
\midrule
\textbf{Subtotal Ψ-Modulation} & 44 \\
\midrule
\textbf{TOTAL} & \textbf{188} \\
\bottomrule
\end{tabular}
\end{table}

%------------------------------------------------------------------------------
\subsubsection{Implications for the 9 Application Cases}
%------------------------------------------------------------------------------

The Ψ-modulation explains \textbf{treatment effect heterogeneity} across contexts:

\begin{table}[htbp]
\centering
\caption{Context-Dependence of Application Cases}
\label{tab:case-psi-dependence}
\small
\renewcommand{\arraystretch}{1.2}
\begin{tabular}{@{}llll@{}}
\toprule
\textbf{Case} & \textbf{Key Ψ-Dims} & \textbf{High-Context Prediction} & \textbf{Low-Context Prediction} \\
\midrule
CASE-01 Labor & $\Psi_S$, $\Psi_F$ & Easier identity bridges & Harder transitions \\
CASE-02 Smoking & $\Psi_E$, $\Psi_K$ & Better info penetrates & Filter blocks info \\
CASE-03 Saving & $\Psi_E$, $\Psi_T$ & Longer planning horizon & Extreme myopia \\
CASE-04 Climate & $\Psi_S$, $\Psi_I$ & KNU more accessible & KNU invisible \\
CASE-05 Education & $\Psi_E$, $\Psi_C$ & SES effect weaker & SES dominates \\
CASE-06 Polarization & $\Psi_K$, $\Psi_I$ & Weaker filter & Maximal filtering \\
CASE-07 Tech & $\Psi_F$, $\Psi_M$ & Faster adoption & Structural barriers \\
CASE-08 Org Change & $\Psi_S$, $\Psi_I$ & Cascades manageable & Cascades dominate \\
CASE-09 Sustainable & $\Psi_S$, $\Psi_K$ & Smaller gap & Larger gap \\
\bottomrule
\end{tabular}
\end{table}

\begin{cross-reference}
\textbf{Complete Specification:} Full Ψ-dimension definitions, all interaction coefficients, and calibration methodology are documented in Appendix~\ref{app:bcm-axiom-formalization}, Section ``Ψ-Context Modulation of Awareness.''
\end{cross-reference}

%%%%%%%%%%%%%%%%%%%%%%%%%%%%%%%%%%%%%%%%%%%%%%%%%%%%%%%%%%%%%%%%%%%%%%%%%%%%%%%
% END OF SECTION 11.15
%%%%%%%%%%%%%%%%%%%%%%%%%%%%%%%%%%%%%%%%%%%%%%%%%%%%%%%%%%%%%%%%%%%%%%%%%%%%%%%
